\section{Related Work}
\label{sec:related}

\subsection{Spurious Correlations and Worst-Group Accuracy}

Spurious correlations --- statistical associations between input features and labels that hold in training data but not at test time --- are a well-documented failure mode of empirical risk minimization \cite{sagawa2019distributionally}. On the Waterbirds benchmark \cite{sagawa2019distributionally}, ResNet-50 models trained with ERM exploit the background attribute (land vs.\ water) as a shortcut for bird species classification, achieving high average accuracy but failing on minority groups (e.g., landbirds on water). WGA --- the accuracy on the worst-performing group --- has become the standard measure for spurious correlation robustness.

Multiple training-time interventions improve WGA over ERM: GroupDRO \cite{sagawa2019distributionally} reweights groups dynamically; JTT \cite{liu2021just} upweights misclassified examples from a warm-started ERM model; SAM \cite{foret2021sharpness} improves sharpness-aware generalization. All achieve WGA improvements on Waterbirds, yet \emph{why} they improve --- specifically, what mechanism drives the improvement --- remains underspecified. Our study addresses this mechanistic gap for GroupDRO and DFR.

\subsection{Deep Feature Reweighting and Last-Layer Retraining}

Kirichenko et al.\ \cite{kirichenko2022last} showed that much of WGA improvement can be achieved by simply retraining the final classification head on a group-balanced held-out set, while freezing the backbone entirely (DFR). Subsequent work \cite{labonte2023last,hill2025unreasonable} refined this understanding: the key condition is group balance in the held-out set used for head retraining. Importantly, DFR and ERM backbones should be \emph{identical at the weight level} by design. Izmailov et al.\ \cite{izmailov2022feature} release both DFR and ERM checkpoints but do not explicitly quantify this backbone identity empirically. We fill this gap in H-P0 (cosine similarity = 1.000000 $\pm$ 1e-14, all 3 seed pairs).

Le et al.\ \cite{le2023last} provide a cautionary counterpoint: in medical imaging domains, DFR is insufficient when backbone features strongly encode spurious attributes, and backbone-level intervention is required. This corroborates our framework: the sufficiency of head-only retraining depends on the backbone's spurious encoding level.

\subsection{Feature Analysis of Robustification Methods}

Izmailov et al.\ \cite{izmailov2022feature} use a spurious-DFR proxy (s-DFR) to estimate how much information ERM backbones encode about spurious attributes ($\sim$85\% probe accuracy), but do not report per-method, per-seed probe accuracy with statistical tests comparing GroupDRO to ERM. Murotkar et al.\ \cite{murotkar2024identifying} use sklearn L-BFGS linear probes on frozen ResNet-50 layer4 features for spurious feature disentanglement --- establishing the probe methodology we adopt.

Park et al.\ \cite{park2025spurious} introduce SCER, which explicitly regularizes the spurious feature subspace during training to reduce background decodability. SCER demonstrates a causal link between reducing background decodability and improving WGA --- a link our study validates from the opposite direction: GroupDRO \emph{implicitly} achieves what SCER does explicitly. Raymond et al.\ \cite{raymond2026groupdro} provide corroborating evidence that GroupDRO reshapes backbone representations across all layers, consistent with our H-M2 findings; our H-M2 result stands on its own pre-registered weight-difference evidence independent of this citation.

\subsection{Linear Probing for Representation Analysis}

Linear probing is a standard tool for analyzing backbone representation content \cite{alain2016understanding}. The conventions we follow (sklearn L-BFGS, $C=10^9$, frozen layer4 features with global average pooling) match those established by Kirichenko et al.\ \cite{kirichenko2022last} and Murotkar et al.\ \cite{murotkar2024identifying}. Our contribution is the \emph{measurement} of backbone spurious encoding across methods with pre-registered statistical tests and per-seed granularity --- enabling the backbone-vs-head typology.
